% Resultados
\section{Resultados}

\subsection{Eje 1: Distribución de la Pobreza}

\subsubsection{Incidencia de pobreza por región}

\begin{table}[H]
    \centering
    \caption{Incidencia de pobreza por región (CASEN 2022, diseño complejo)}
    \label{tab:pobreza_region}
    \begin{tabular}{lccc}
    \toprule
    \tableheadercolor
    \textbf{Región} & \textbf{Prop.} & \textbf{IC 95\%} & \textbf{n} \\
    \midrule
    La Araucanía & 9.97\% & [8.95\%, 11.0\%] & 4,916 \\
    Ñuble & 9.97\% & [8.78\%, 11.2\%] & 3,308 \\
    Tarapacá & 9.36\% & [7.23\%, 11.5\%] & 2,806 \\
    Atacama & 7.55\% & [6.17\%, 8.93\%] & 3,211 \\
    Arica y Parinacota & 7.42\% & [6.16\%, 8.69\%] & 2,718 \\
    Maule & 7.36\% & [6.56\%, 8.16\%] & 5,169 \\
    Biobío & 7.15\% & [6.46\%, 7.84\%] & 7,076 \\
    \midrule
    Metropolitana & 3.73\% & [3.35\%, 4.11\%] & 13,202 \\
    Aysén & 3.38\% & [2.26\%, 4.49\%] & 1,449 \\
    Magallanes & 2.73\% & [1.55\%, 3.91\%] & 1,964 \\
    \bottomrule
    \end{tabular}
    
    \vspace{2pt}
    {\small\textit{Nota: Estimaciones con diseño muestral complejo (estratos y conglomerados). Se muestran las 7 regiones con mayor y las 3 con menor incidencia.}}
\end{table}

Existe una marcada heterogeneidad territorial. Las regiones de La Araucanía y Ñuble presentan las mayores tasas de pobreza (9.97\% cada una, IC 95\%: $\sim$9--11\%), seguidas por Tarapacá (9.36\%) y Atacama (7.55\%). En contraste, las regiones extremas de Magallanes (2.73\%, IC: [1.55\%, 3.91\%]) y Aysén (3.38\%), junto con la Región Metropolitana (3.73\%), presentan los menores niveles de pobreza entre jefes de hogar. La Figura \ref{fig:pobreza_region} ilustra esta distribución.

\begin{figure}[H]
    \centering
    \begin{subfigure}[b]{0.48\textwidth}
        \centering
        \includegraphics[width=\textwidth]{Imagenes/g1_pobreza_region.png}
        \caption{Gráfico: Pobreza por región}
        \label{fig:pobreza_region_graf}
    \end{subfigure}
    \hfill
    \begin{subfigure}[b]{0.38\textwidth}
        \centering
        \includegraphics[width=\textwidth]{Imagenes/mapa_chile.pdf}
        \caption{Mapa de Chile}
        \label{fig:mapa_chile}
    \end{subfigure}
    \caption{Incidencia de pobreza por región (CASEN 2022)}
    \label{fig:pobreza_region}
\end{figure}

\subsubsection{Pobreza por zona (urbano/rural)}

\begin{table}[H]
    \centering
    \caption{Pobreza por zona urbana y rural (diseño complejo)}
    \label{tab:pobreza_zona}
    \begin{tabular}{lcccc}
    \toprule
    \tableheadercolor
    \textbf{Zona} & \textbf{Proporción} & \textbf{E.E.} & \textbf{IC 95\%} & \textbf{n expandido} \\
    \midrule
    Urbano & 5.28\% & 0.12\% & [5.04\%, 5.52\%] & 6,193,975 \\
    Rural & 8.37\% & 0.30\% & [7.79\%, 8.95\%] & 804,118 \\
    \bottomrule
    \end{tabular}
    
    \vspace{2pt}
    {\small\textit{Nota: Diferencia de 3.09 pp, significativa ($\chi^2$ Rao-Scott = 131.79, p $<$ 0.001).}}
\end{table}

La pobreza rural es significativamente mayor que la urbana, con una diferencia relativa del 58\%. El test Chi-cuadrado de Rao \& Scott confirma la asociación ($\chi^2 = 131.79$, $df = 1$, $p < 2.2 \times 10^{-16}$), y el test t de diseño complejo arroja una diferencia de 3.09 puntos porcentuales ($t = 9.62$, $df = 11306$, $p < 0.001$, IC 95\%: [2.46\%, 3.72\%]). La Figura \ref{fig:pobreza_urbano_rural} visualiza esta brecha.

\begin{figure}[H]
    \centering
    \begin{subfigure}[b]{0.48\textwidth}
        \centering
        \includegraphics[width=\textwidth]{Imagenes/g4_pobreza_urbano_rural.png}
        \caption{Pobreza urbano vs rural}
        \label{fig:pobreza_urbano_rural}
    \end{subfigure}
    \hfill
    \begin{subfigure}[b]{0.48\textwidth}
        \centering
        \includegraphics[width=\textwidth]{Imagenes/g7_pobreza_region_zona.png}
        \caption{Pobreza por región y zona}
        \label{fig:pobreza_region_zona}
    \end{subfigure}
    \caption{Análisis de pobreza territorial}
    \label{fig:pobreza_territorial}
\end{figure}

\subsubsection{Modelo logístico de pobreza}

\begin{table}[H]
    \centering
    \caption{Estimaciones del modelo logístico de pobreza}
    \label{tab:modelo_logistico}
    \begin{tabular}{lccp{4cm}}
    \toprule
    \tableheadercolor
    \textbf{Variable} & \textbf{OR} & \textbf{IC 95\%} & \textbf{Interpretación} \\
    \midrule
    Zona rural (ref: urbano) & 1.34 & [1.22, 1.47] & +34\% odds de pobreza \\
    Escolaridad (por año) & 0.93 & [0.93, 0.94] & $-$7\% odds por año \\
    Tamaño hogar (por persona) & 1.21 & [1.19, 1.24] & +21\% odds por persona \\
    \bottomrule
    \end{tabular}
    
    \vspace{2pt}
    {\small\textit{Nota: Modelo quasibinomial con diseño complejo. Dispersión = 0.98. Todas las variables $p < 2 \times 10^{-16}$.}}
\end{table}

\subsubsection{Ingreso por tamaño de hogar}

\begin{table}[H]
    \centering
    \caption{Ingreso total del hogar por tamaño (CASEN 2022)}
    \label{tab:ingreso_tamano}
    \begin{tabular}{lccc}
    \toprule
    \tableheadercolor
    \textbf{Tamaño hogar} & \textbf{Ingreso medio} & \textbf{IC 95\%} \\
    \midrule
    1 persona & \$1,061,844 & [\$1,024,718; \$1,098,970] \\
    2 personas & \$1,462,260 & [\$1,411,060; \$1,513,460] \\
    3--4 personas & \$1,738,278 & [\$1,703,459; \$1,773,097] \\
    5+ personas & \$2,006,078 & [\$1,942,601; \$2,069,556] \\
    \bottomrule
    \end{tabular}
    
    \vspace{2pt}
    {\small\textit{Nota: Variable \texttt{ytotcorh} (ingreso total corregido del hogar).}}
\end{table}

La educación es el factor protector más potente. Cada año adicional de escolaridad reduce el odds de pobreza en \textbf{7\%}, controlando por zona y tamaño del hogar. El modelo presenta un parámetro de dispersión de 0.98, indicando buen ajuste. Todas las variables son altamente significativas ($p < 2 \times 10^{-16}$). La Figura \ref{fig:forest_plot} presenta los odds ratios con sus intervalos de confianza.

\begin{figure}[H]
    \centering
    \includegraphics[width=0.55\textwidth]{Imagenes/g9_forest_plot_pobreza.png}
    \caption{Forest plot de odds ratios del modelo logístico}
    \label{fig:forest_plot}
\end{figure}

\subsubsection{Análisis de mediación}

\begin{table}[H]
    \centering
    \caption{Descomposición de efectos: Zona $\to$ Educación $\to$ Pobreza}
    \label{tab:mediacion}
    \begin{tabular}{lcp{4.5cm}}
    \toprule
    \tableheadercolor
    \textbf{Efecto} & \textbf{Valor} & \textbf{Interpretación} \\
    \midrule
    Efecto total (c) & 0.4949 & Log-odds pobreza rural vs urbano \\
    Efecto directo (c') & 0.2746 & Efecto neto controlando educación \\
    Efecto indirecto (a$\times$b) & 0.2224 & Vía menor educación rural \\
    \midrule
    \textbf{Proporción mediada} & \textbf{44.9\%} & Educación explica casi la mitad \\
    Test de Sobel & Z=16.795, p$<$0.001 & Mediación significativa \\
    \bottomrule
    \end{tabular}
\end{table}

La menor escolaridad en zonas rurales media el 44.9\% del efecto de la ruralidad sobre la pobreza (Figura \ref{fig:mediacion}). El análisis de mediación se realizó con diseño muestral complejo, filtrando observaciones con valores faltantes. Políticas de acceso educativo rural podrían reducir sustancialmente esta brecha.

\begin{figure}[H]
    \centering
    \begin{subfigure}[b]{0.48\textwidth}
        \centering
        \includegraphics[width=\textwidth]{Imagenes/g10_diagrama_mediacion.png}
        \caption{Diagrama de mediación}
        \label{fig:mediacion}
    \end{subfigure}
    \hfill
    \begin{subfigure}[b]{0.48\textwidth}
        \centering
        \includegraphics[width=\textwidth]{Imagenes/g6_escolaridad_zona.png}
        \caption{Escolaridad por zona}
        \label{fig:escolaridad_zona}
    \end{subfigure}
    \caption{Análisis de mediación: Zona $\to$ Educación $\to$ Pobreza}
    \label{fig:mediacion_combined}
\end{figure}

%-------------------------------------------------------------------------------
\subsection{Eje 2: Brecha Salarial de Género}
%-------------------------------------------------------------------------------

\subsubsection{Brecha bruta (sin controles)}

\begin{table}[H]
    \centering
    \caption{Ingresos laborales por sexo (brecha bruta)}
    \label{tab:ingreso_bruto}
    \begin{tabular}{lccc}
    \toprule
    \tableheadercolor
    \textbf{Indicador} & \textbf{Hombre} & \textbf{Mujer} & \textbf{Brecha} \\
    \midrule
    Ingreso medio & \$1,402,711 & \$1,112,358 & \textbf{20.7\%} \\
    IC 95\% (media) & [\$1,365,373; \$1,440,050] & [\$1,087,319; \$1,137,398] & --- \\
    Mediana & \$920,000 & \$780,000 & 15.2\% \\
    \bottomrule
    \end{tabular}
    
    \vspace{2pt}
    {\small\textit{Nota: Test t de diseño complejo: $t = -14.53$, $df = 11261$, diferencia = \$290,353 ($p < 2.2 \times 10^{-16}$).}}
\end{table}

Las jefas de hogar perciben en promedio un 20.7\% menos de ingresos laborales. La brecha en medianas (15.2\%) es menor, sugiriendo que la cola superior de la distribución masculina amplifica la diferencia en medias (Figura \ref{fig:distribucion_ingreso}).

\begin{figure}[H]
    \centering
    \includegraphics[width=0.55\textwidth]{Imagenes/g5_distribucion_ingreso_sexo.png}
    \caption{Distribución de ingresos laborales por sexo}
    \label{fig:distribucion_ingreso}
\end{figure}

\subsubsection{Brecha por nivel educativo}

\begin{table}[H]
    \centering
    \caption{Brecha salarial por nivel educativo}
    \label{tab:brecha_educacion}
    \begin{tabular}{lcp{4.5cm}}
    \toprule
    \tableheadercolor
    \textbf{Nivel educativo} & \textbf{Brecha} & \textbf{Interpretación} \\
    \midrule
    Básica o menos & \textbf{8.2\%} & Menor brecha relativa \\
    Media & 16.9\% & Cerca del promedio nacional \\
    Técnico Superior & \textbf{21.1\%} & Mayor brecha absoluta \\
    Profesional & 14.1\% & Menor que técnico superior \\
    Postgrado & \textbf{20.8\%} & Segunda mayor brecha \\
    \bottomrule
    \end{tabular}
\end{table}

\textbf{Hallazgo clave:} La brecha salarial presenta un patrón \textbf{no lineal} con la educación (Figura \ref{fig:brecha_educacion_combined}). Es mínima en educación básica (8.2\%), aumenta drásticamente en media y técnica superior (16.9\%--21.1\%), disminuye en profesionales (14.1\%), y vuelve a subir en postgrado (20.8\%). Este patrón sugiere que la brecha no se explica únicamente por capital humano.

% Figuras lado a lado: Ingreso por educación y Brecha por educación
\begin{figure}[H]
    \centering
    \begin{subfigure}[b]{0.48\textwidth}
        \centering
        \includegraphics[width=\textwidth]{Imagenes/g2_ingreso_educacion_sexo.png}
        \caption{Ingreso por nivel educativo y sexo}
        \label{fig:ingreso_educacion}
    \end{subfigure}
    \hfill
    \begin{subfigure}[b]{0.48\textwidth}
        \centering
        \includegraphics[width=\textwidth]{Imagenes/g3_brecha_educacion.png}
        \caption{Brecha salarial porcentual}
        \label{fig:brecha_educacion_graf}
    \end{subfigure}
    \caption{Análisis de brecha salarial por nivel educativo}
    \label{fig:brecha_educacion_combined}
\end{figure}

\subsubsection{Modelos de regresión}

\begin{table}[H]
    \centering
    \caption{Modelos de brecha salarial ajustados}
    \label{tab:modelos_brecha}
    \begin{tabular}{lcc}
    \toprule
    \tableheadercolor
    \textbf{Modelo} & \textbf{Coef. Sexo (Mujer)} & \textbf{Interpretación} \\
    \midrule
    Ajustado (educación, edad) & $-$\$243,330 (E.E.=\$18,112) & $t = -13.44$, $p < 0.001$ \\
    Log-lineal (educación, edad, edad$^2$) & $-$0.200 (E.E.=0.010) & $-$18.2\% ingreso \\
    Completo (+ ocupación, hogar) & $-$\$169,923 (E.E.=\$26,000) & $p = 8.8 \times 10^{-11}$ \\
    \bottomrule
    \end{tabular}
    
    \vspace{2pt}
    {\small\textit{Nota: Todos los modelos estimados con diseño muestral complejo. Se rechaza H$_0$ en todos los casos.}}
\end{table}

Controlando por educación, edad, ocupación y tamaño del hogar, las jefas de hogar perciben un \textbf{18.2\%} menos de ingreso laboral (modelo log-lineal). Se \textbf{rechaza H$_0$}: existe brecha salarial significativa contra las mujeres incluso con controles completos.

\subsubsection{Brecha por grupo de edad}

\begin{table}[H]
    \centering
    \caption{Brecha salarial por grupo de edad}
    \label{tab:brecha_edad}
    \begin{tabular}{lcp{5cm}}
    \toprule
    \tableheadercolor
    \textbf{Grupo de edad} & \textbf{Brecha} & \textbf{Interpretación} \\
    \midrule
    $\leq$30 años & \textbf{13.1\%} & Menor brecha (entrada al mercado) \\
    31--40 años & 20.8\% & Cerca del promedio ($\sim$20\%) \\
    41--50 años & \textbf{24.5\%} & Alta brecha (pico de carrera) \\
    51--60 años & 17.8\% & Reducción respecto al anterior \\
    60+ años & \textbf{24.8\%} & Mayor brecha (pre-jubilación) \\
    \bottomrule
    \end{tabular}
\end{table}

La brecha salarial presenta un patrón de \textbf{U invertida con rebote} (Figura \ref{fig:brecha_edad}). Es menor al inicio de la carrera laboral (13.1\%), alcanza su máximo en 41--50 años (24.5\%), disminuye en 51--60 (17.8\%), pero vuelve a aumentar en mayores de 60 (24.8\%). Esto sugiere efectos de cohorte y posible discriminación en etapas avanzadas.

\begin{figure}[H]
    \centering
    \includegraphics[width=0.52\textwidth]{Imagenes/g8_brecha_edad.png}
    \caption{Brecha salarial de género por grupo de edad}
    \label{fig:brecha_edad}
\end{figure}

%-------------------------------------------------------------------------------
\subsection{Validación: Sesgo por No Ponderar}
%-------------------------------------------------------------------------------

\begin{table}[H]
    \centering
    \caption{Sesgo por ignorar factores de expansión}
    \label{tab:sesgo_ponderacion}
    \begin{tabular}{lc}
    \toprule
    \tableheadercolor
    \textbf{Estimador} & \textbf{Valor} \\
    \midrule
    Media NO ponderada & \$1,321,458 \\
    Media ponderada & \$1,566,277 \\
    \textbf{Diferencia relativa} & \textbf{+18.5\%} \\
    \bottomrule
    \end{tabular}
\end{table}

Ignorar los factores de expansión subestima el ingreso promedio en $\sim$18.5\%. La muestra CASEN sobrerrepresenta zonas de mayor densidad poblacional, y los pesos corrigen esta distorsión.

\clearpage
